\section{Modern Video-Compression Standards}
\subsection{Universal Hybrid Video Encoder}
From MPEG-1 going on they use all the same structure, core points:
\begin{itemize}
    \item Spatial (intra) and temporal (inter) prediction
    \item Residual calculation
    \item Transform and Quantization
    \item Internal closed loops to provide reliability
\end{itemize}
Rule of thumb: every generation reduces video size by $1/2$ and increases complexity by $10\times$
\begin{center}
\includegraphics[width=0.95\linewidth]{video_chain.png}
\end{center}
\noindent End-to-end video chain. Only the \textbf{decoding} is normative (\emph{scope of the standard}): pre/post-processing and the encoder are left to implementation competition.
\begin{center}
\includegraphics[width=\linewidth]{modern_encoder.png}
\end{center}
\noindent Modern hybrid encoder (H.264-style): integer two-stage $4\times4$ DCT, in-loop \textbf{Deblocking Filter} inside the reconstruction loop before the Frame Buffer, and \textbf{CABAC/CAVLC} entropy coding. Same closed-loop structure as the universal encoder, refined block by block.
\subsubsection{Codecs applications}
\begin{itemize}
    \item Real Time $\to$ H.264 (very diffused, also in HW implementations since 2003)
    \item Broadcast $\to$ H.264 + HEVC, growing HEVC market share
    \item Web Streaming $\to$ H.264 + HEVC + AV1 (still royalty free)
    \item Next-Generations $\to$ need to use massive bitrates of raw video $\to$ VVC, AV1
\end{itemize}
\subsection{Rate-Distorsion Optimization}
$$J = D + \lambda R$$
Using the R-D optimization method we have a huge number of options, how do we choose the best mode?\\
The solution is to work in blocks $\to$ Block Partitioning Problem
\subsubsection{Block Partitioning Problem: H.264}
In H.264 we use macroblocks $16 \times 16$, it allows further splitting until $4 \times 4$ blocks, useful to compute edges.\\
Still this size is too low for High-Resoluton videos, we need something more.
\subsubsection{Block Partitioning Problem: H.265/HEVC}
HEVC uses a Quad-Tree paradigm, it is like a 2D binary tree.\\
A certain CTU (Code Tree Unit) can be split into 4 CU (Coding Units), then a CU becomes a CTU and we can apply this recursively.
Standards should provide only the syntax, we need to do the actual implementation.\\
We continue splitting until another split provides no more improvement, and complexity becomes reasonable.\\
CU can be split even in asymmetrical quantities into the PU (Prediction Units), then we use S and T prediction.
\subsubsection{Block Partitioning Problem: H.266/VVC}
It uses the same paradigm but here we have a 10-way superblock partitioning, we can binary and ternary split.\\
How can we build the decision tree? Depth First Search:
\begin{itemize}
    \item Encode Block, freeze it
    \item Go to the next neighbor block
    \item Iterate until all blocks are encoded
\end{itemize}
If the block is flat, doesn't make any sense to partition it.\\
In this case we also have a huge number of possible partitions, but we give up in looking all possible choices.\\
This is like a classification among pixel groups, ML and NN are good in solving these problems.\\
Researchers are looking for a light NN to exploit this.
\subsection{Spatial and Temporal Frame Prediction}
\subsubsection{Intra Prediction}
We want to exploit spatial prediction among neighbors.
Instead of trying just one prediction we use different predictors in parallel for the same block, then we pick the best one. We have to signal this to the Encoder.\\
We take one director-per-block instead of one-per-pixel (it is too much bits wasted):
% Integrated from Summaries: exact mode counts per standard
\begin{itemize}
    \item \textbf{H.264/AVC}: 9 modes for $4\times4$ blocks (\textbf{8 directional + DC}); 4 modes for $16\times16$ blocks (2 bits). A $16\times16$ MB split in 16 $4\times4$ subblocks has ideally $9^{16}$ combinations: done sequentially to cut them down.
    \item \textbf{H.265/HEVC}: \textbf{35 modes = 33 directional + DC + Planar}, to handle complex textures.
    \item \textbf{H.266/VVC}: \textbf{65 directional modes + wide-angle} modes for non-square blocks.
\end{itemize}
Predictors grow exponentially with the signaling bits: $P \propto 2^B$. With 67 modes, $\approx 7$ bits/block of pure signaling would destroy the gain on $4\times4$ blocks.\\
\textbf{MPM: Most Probable Mode} idea, ENC predicts which is the most probable direction using the spatial context: a candidate list is built from the modes chosen by top and left neighbors (HEVC: 3 candidates, VVC: 6). If the chosen mode is in the list, 1--2 bits suffice; only if it fails long codewords are used.
\subsubsection{Inter Prediction}
Here we use Fractional Motion Compensation, each frame is computed using multiple reference frames.\\
Foreach image in the GOP, we have to say what is the List of frames used for prediction.\\
 Here complexity $\propto$ size of bits. How do we encode the motion vectors?\\
A better strategy is to produce a list of predictors and encode the indexes.
\subsubsection{Merge Mode}
Used in AV1/VVC, not only picks the pixels but also applies affine transformation on the blocks, thus producing a complexity increase.
\subsection{Filtering, Transforms and Quantization}
\subsubsection{Residual Coding}
DCT problem: we want floating number operations (too complex for an encoder, must be quick).\\
Solution: just use integer power of 2, such that multiplications are bit shifts and we have also $\pm$ operations.\\
We don't provide quantization table, but a similar Quantization Parameter ($\propto$ dB step size)\\
% Integrated from Summaries
\textbf{QP behavior} (H.264/HEVC): the quantization step \textbf{doubles every +6 of QP}; empirically +1 QP $\approx$ $-12.5\%$ rate. QP is the main knob for rate control.\\
Entropy coding standardized around \textbf{CABAC} (Context-Adaptive Binary Arithmetic Coding): binarizes syntax elements and adapts probabilities on spatial context; $5\to15\%$ rate saving over older VLC (CAVLC).
% Integrated from Summaries: deblocking is a sample exam question
\subsubsection{In-loop Deblocking Filter}
\textbf{Why artifacts}: transform and quantization are applied \emph{independently on disjoint blocks} $\Rightarrow$ at low bitrate strong discontinuities (\textbf{blocking artifacts}) appear at block boundaries.\\
\textbf{Why in-loop} (and not just post-processing): uncorrected artifacts would be stored in the Decoded Picture Buffer and copied around by future motion compensations, propagating the damage. The filter is \textbf{normative} (inside the reconstruction loop, identical at encoder and decoder to keep sync).\\
\textbf{Key insight}: artifact locations are known a priori --- the block edges --- so we apply adaptive smart smoothing there, without blurring real image edges.
\begin{itemize}
    \item \textbf{H.264/AVC deblocking}: analyzes edges between $4\times4$ blocks; filtering strength adapts dynamically to coding mode, motion vectors and quantization step of the neighboring blocks.
    \item \textbf{H.265/HEVC}: deblocking only on $8\times8$ grid (less complexity) + \textbf{SAO} (Sample Adaptive Offset): classifies pixels by edges/bands and adds offsets to fix ringing artifacts.
    \item \textbf{H.266/VVC}: adds \textbf{ALF} (Adaptive Loop Filter) and LMCS (HDR detail preservation).
    \item \textbf{AV1}: CDEF directional filters, Wiener loop-restoration filters, film grain synthesis (grain re-synthesized at decoder at zero rate cost).
\end{itemize}
\subsubsection{Slices}
Error corruption behavior: artifacts highlight due to error propagation in compression.\\
How can we prevent this? Slices are subparts of the frame, that are independently encoded, such that we can parallel-encode them, because we don't use them for prediction.\\
At Intra Frames CABAC\footnote{Context-Adaptive Binary Arithmetic Coding} is reset.
\subsection{Network parallelism}
All the video data is encoded into a bitstream.
\subsubsection{NALU}
\begin{itemize}
    \item VCL NALUs: contain all the compression data for a slice, can be IDR\footnote{Istantaneous Decoding Refresh, video access point} (I + CABAC reset) or non-IDR
    \item non-VCL NALUs: contain SPS, PPS, SEI
    \begin{itemize}
        \item SPS: Sequence Parameter Set (picture size, color depth,...)
        \item PPS: Picture Parameter Set (coding specific parameters)
        \item SEI: Supplemenetal and Enhancement Infos (Subtitles, HDR infos, ...)
    \end{itemize}
\end{itemize}
\begin{center}
\includegraphics[width=0.45\linewidth]{h264_layers.png}
\end{center}
\noindent H.264 layered architecture: the \textbf{Video Coding Layer} produces the compressed bitstream, the \textbf{Network Abstraction Layer} wraps it into NAL units, which the Transport Layer packetizes for the network.
NALUs are then packetized into a sequential bit stream, there are 2 ways to signal the NALU (H.264/HEVC):
\begin{itemize}
    \item Start Code $0x00000001$: it is used but we have problems if this sequence is repeated in a nalu (start code mismatch), we have to use Emulation Prevention Techniques
    \item Length Header: before the NALU we specify the actual NALU length, such that we can read it or skip it.
\end{itemize}
Actual Sequence using Annex-B ordering is:
$$\begin{aligned}
    \boxed{\text{SPS}}\boxed{\text{PPS}}&\boxed{\text{AUD}}\boxed{\text{SEI}}\boxed{\text{I slice }1\to 4}\\
    &\boxed{\text{AUD}}\boxed{\text{SEI}}\boxed{\text{P slice }1\to 2}\\
    &\boxed{\text{AUD}}\boxed{\text{SEI}}\boxed{B_1\text{I slice }1}\\
    &\boxed{\text{AUD}}\boxed{\text{SEI}}\boxed{B_2\text{ slice }1}\\
\end{aligned}
$$
AV1 drops this approach and uses OBUs (Open Bitstream Units), in CDNs there are MANE\footnote{Media Aware Network Elements} that can allow to change the bitstream on the fly to add redundance or remove unnecessary NALUs.\\
\begin{center}
\includegraphics[width=0.9\linewidth]{nal_stream.png}
\end{center}
\noindent NAL-unit stream Encoder $\to$ MANE $\to$ Decoder: parameter sets (SPS, PPS), SEI and slice NALUs are forwarded packet by packet.
\begin{center}
\includegraphics[width=0.9\linewidth]{nal_prune.png}
\end{center}
\noindent \textbf{Pruning}: the MANE drops the non-decodable B-frame NALUs, forwarding only the I and P frames (rate adaptation without re-encoding).
\begin{center}
\includegraphics[width=0.9\linewidth]{nal_crc.png}
\end{center}
\noindent \textbf{Adding redundancy}: the MANE can inject extra protection (e.g.\ a CRC NALU) into the stream for more robust delivery.
Decoder accesses video at RAPs (Random Access Points), there are 2 main access techniques:
\begin{itemize}
    \item IDR: Technique to clear previous buffer at a certain intra frame. It breaks bidirectional temp. prediction
    \item CRA (Clear Random Access): we just skip some images because they are not decodable.
\end{itemize}